\begin{thebibliography}{99}
\small\raggedright
\bibitem{AgrawalVinay}
Manindra Agrawal and V. Vinay. Arithmetic circuits: A chasm at depth four. In \emph{49th Annual
IEEE Symposium on Foundations of Computer Science}, pages 67--75, 2008.

\bibitem{AKV20}
Noga Alon, Mrinal Kumar, and Ben Lee Volk. Unbalancing sets and an almost quadratic lower
bound for syntactically multilinear arithmetic circuits. \emph{Combinatorica}, 40(2):149--178, 2020.

\bibitem{Batcher}
K. E. Batcher. Sorting networks and their applications. In \emph{Proceedings of the Spring Joint
Computer Conference}, volume 32 of \emph{AFIPS Conference Proceedings}, pages 307--314, 1968.

\bibitem{BaurStrassen}
Walter Baur and Volker Strassen. The complexity of partial derivatives. \emph{Theoretical Computer
Science}, 22(3):317--330, 1983.

\bibitem{BCGTV}
Eli Ben-Sasson, Alessandro Chiesa, Daniel Genkin, Eran Tromer, and Madars Virza. SNARKs
for C: Verifying program executions succinctly and in zero knowledge. In \emph{CRYPTO}, pages
90--108, 2013.

\bibitem{BDS24}
C. S. Bhargav, Sagnik Dutta, and Nitin Saxena. Improved lower bound, and proof barrier, for
constant depth algebraic circuits. \emph{ACM Transactions on Computation Theory}, 16(4):23:1--23:22,
2024.

\bibitem{be}
Markus Bläser and Christian Engels. Parameterized Valiant's classes. In \emph{14th International
Symposium on Parameterized and Exact Computation}, volume 148 of \emph{Leibniz International
Proceedings in Informatics}, pages 3:1--3:14. Schloss Dagstuhl--Leibniz-Zentrum für Informatik,
2019.

\bibitem{BGS}
Mark Braverman, Sumegha Garg, and Ariel Schvartzman. Coding in undirected graphs is either
very helpful or not helpful at all. In \emph{8th Innovations in Theoretical Computer Science Conference
(ITCS 2017)}, volume 67 of \emph{Leibniz International Proceedings in Informatics}, pages 18:1--18:18,
2017.

\bibitem{Bridy}
Andrew Bridy. Automatic sequences and curves over finite fields. \emph{Algebra \& Number Theory},
11(3):685--712, 2017.

\bibitem{Burgisser}
Peter Bürgisser. \emph{Completeness and Reduction in Algebraic Complexity Theory}, volume 7 of
\emph{Algorithms and Computation in Mathematics}. Springer, 2000.

\bibitem{ApproxFactors}
Peter Bürgisser. The complexity of factors of multivariate polynomials. \emph{Foundations of
Computational Mathematics}, 4(4):369--396, 2004. Corrected version: arXiv:1812.06828v1.

\bibitem{BC2024}
Peter Bürgisser. Completeness classes in algebraic complexity theory, 2024. arXiv:2406.06217v1.

\bibitem{BCS}
Peter Bürgisser, Michael Clausen, and Mohammad Amin Shokrollahi. \emph{Algebraic Complexity
Theory}, volume 315 of \emph{Grundlehren der mathematischen Wissenschaften}. Springer, 1997.

\bibitem{BIPnooccurrence}
Peter Bürgisser, Christian Ikenmeyer, and Greta Panova. No occurrence obstructions in geometric
complexity theory. \emph{Journal of the American Mathematical Society}, 32(1):163--193, 2019.

\bibitem{BLMW}
Peter Bürgisser, J. M. Landsberg, Laurent Manivel, and Jerzy Weyman. An overview of
mathematical issues arising in the geometric complexity theory approach to
$\mathrm{VP}\ne \mathrm{VNP}$. \emph{SIAM
Journal on Computing}, 40:1179--1209, 2011.

\bibitem{CaiChenLi}
Jin-Yi Cai, Xi Chen, and Dong Li. Quadratic lower bound for permanent vs. determinant in
any characteristic. \emph{Computational Complexity}, 19:37--56, 2010.

\bibitem{Cartier}
Pierre Cartier. Une nouvelle opération sur les formes différentielles. \emph{Comptes rendus de
l'Académie des sciences de Paris}, 244:426--428, 1957.

\bibitem{CKSS24}
Prerona Chatterjee, Deepanshu Kush, Shubhangi Saraf, and Amir Shpilka. Lower bounds for
set-multilinear branching programs. In \emph{39th Computational Complexity Conference}, volume 300
of \emph{LIPIcs}, pages 20:1--20:20, 2024.

\bibitem{CKW}
Xi Chen, Neeraj Kayal, and Avi Wigderson. Partial derivatives in arithmetic complexity and
beyond. \emph{Foundations and Trends in Theoretical Computer Science}, 6(1--2):1--138, 2011.

\bibitem{Anthropic3SAT}
Claude (Anthropic). A $2^{o(n)}$ Time Algorithm for 3-SAT. Hexagon:2610.00577v1, 2026.

\bibitem{Cook71}
Stephen A. Cook. The complexity of theorem-proving procedures. In \emph{Proceedings of the Third
Annual ACM Symposium on Theory of Computing}, pages 151--158. Association for Computing
Machinery, 1971.

\bibitem{DawarWilsenach}
Anuj Dawar and Gregory Wilsenach. Symmetric arithmetic circuits. \emph{Theory of Computing},
21(14):1--32, 2025.

\bibitem{BFLY26}
Théo Borém Fabris, Nutan Limaye, Srikanth Srinivasan, and Amir Yehudayoff. Multilinear
algebraic branching programs and the min-partition rank method. In \emph{41st Computational
Complexity Conference}, volume 383 of \emph{LIPIcs}, pages 22:1--22:20, 2026.

\bibitem{Forbes2024}
Michael A. Forbes. Low-depth algebraic circuit lower bounds over any field. In \emph{39th
Computational Complexity Conference (CCC 2024)}, volume 300 of \emph{Leibniz International
Proceedings in
Informatics}, pages 31:1--31:16. Schloss Dagstuhl -- Leibniz-Zentrum für Informatik, 2024.

\bibitem{FSV18}
Michael A. Forbes, Amir Shpilka, and Ben Lee Volk. Succinct hitting sets and barriers to
proving lower bounds for algebraic circuits. \emph{Theory of Computing}, 14(18):1--45, 2018.

\bibitem{Gro15}
Joshua A. Grochow. Unifying known lower bounds via geometric complexity theory. \emph{Computational
Complexity}, 24:393--475, 2015.

\bibitem{GKSS17}
Joshua A. Grochow, Mrinal Kumar, Michael Saks, and Shubhangi Saraf. Towards an algebraic
natural proofs barrier via polynomial identity testing. Technical Report TR17-009, Electronic
Colloquium on Computational Complexity, 2017.

\bibitem{GSS}
Zeyu Guo, Nitin Saxena, and Amit Sinhababu. Algebraic dependencies and PSPACE algorithms
in approximative complexity over any field. \emph{Theory of Computing}, 15(16):1--30, 2019.

\bibitem{GKKS14}
Ankit Gupta, Pritish Kamath, Neeraj Kayal, and Ramprasad Saptharishi. Approaching the
chasm at depth four. \emph{Journal of the ACM}, 61(6):33:1--33:16, 2014.

\bibitem{Hartshorne}
Robin Hartshorne. \emph{Algebraic Geometry}, volume 52 of \emph{Graduate Texts in Mathematics}. Springer,
1977.

\bibitem{Heintz}
Joos Heintz. Definability and fast quantifier elimination in algebraically closed fields. \emph{Theoretical
Computer Science}, 24:239--277, 1983.

\bibitem{HLW}
Shlomo Hoory, Nathan Linial, and Avi Wigderson. Expander graphs and their applications.
\emph{Bull. Amer. Math. Soc.}, 43:439--561, 2006.

\bibitem{Hrubes}
Pavel Hrube\v{s}. A subquadratic upper bound on Hurwitz's problem and related noncommutative
polynomials. \emph{SIAM Journal on Computing}, 55(2):451--467, 2026.

\bibitem{HWY}
Pavel Hrube\v{s}, Avi Wigderson, and Amir Yehudayoff. Non-commutative circuits and the
sum-of-squares problem. \emph{Journal of the American Mathematical Society}, 24:871--898, 2011.

\bibitem{Hyafil}
Laurent Hyafil. On the parallel evaluation of multivariate polynomials. \emph{SIAM Journal on
Computing}, 8(2):120--123, 1979.

\bibitem{JerrumSnir}
Mark Jerrum and Marc Snir. Some exact complexity results for straight-line computations over
semirings. \emph{Journal of the ACM}, 29(3):874--897, 1982.

\bibitem{Katz}
Nicholas M. Katz. Nilpotent connections and the monodromy theorem: Applications of a result
of Turrittin. \emph{Publications Mathématiques de l'IHÉS}, 39:175--232, 1970.

\bibitem{Kayal2012}
Neeraj Kayal. An exponential lower bound for the sum of powers of bounded degree polynomials.
Technical Report TR12-081, Electronic Colloquium on Computational Complexity, 2012.

\bibitem{KLSS14}
Neeraj Kayal, Nutan Limaye, Chandan Saha, and Srikanth Srinivasan. An exponential lower
bound for homogeneous depth four arithmetic formulas. Technical Report TR14-005, Electronic
Colloquium on Computational Complexity, 2014.

\bibitem{KSS13regular}
Neeraj Kayal, Chandan Saha, and Ramprasad Saptharishi. A super-polynomial lower bound for
regular arithmetic formulas. Technical Report TR13-091, Electronic Colloquium on Computational Complexity, 2013.

\bibitem{KSsurvey}
Neeraj Kayal and Ramprasad Saptharishi. A selection of lower bounds for arithmetic circuits.
In \emph{Perspectives in Computational Complexity}, pages 77--115. Birkhäuser, 2014.

\bibitem{Khoruzhii}
Kirill Khoruzhii. Crossing decomposition for subquadratic noncommutative circuits. Lemmata,
November 2, revised November 4, 2026.

\bibitem{Koiran2012}
Pascal Koiran. Arithmetic circuits: The chasm at depth four gets wider. \emph{Theoretical Computer
Science}, 448:56--65, 2012.

\bibitem{KumarSaraf2014}
Mrinal Kumar and Shubhangi Saraf. On the power of homogeneous depth 4 arithmetic circuits.
\emph{SIAM Journal on Computing}, 46(1):336--387, 2017.

\bibitem{KumarVolk2026}
Mrinal Kumar and Ben Lee Volk. A quadratic lower bound on determinantal complexity.
arXiv:2609.34462v1, 2026.

\bibitem{LandsbergManivelRessayre}
Joseph M. Landsberg, Laurent Manivel, and Nicolas Ressayre. Hypersurfaces with degenerate
duals and the geometric complexity theory program. \emph{Commentarii Mathematici Helvetici},
88(2):469--484, 2013.

\bibitem{LehmkuhlLickteig}
T. Lehmkuhl and T. Lickteig. On the order of approximation in approximative triadic decompositions
of tensors. \emph{Theoretical Computer Science}, 66(1):1--14, 1989.

\bibitem{Levin}
Leonid A. Levin. Universal sequential search problems. \emph{Problems Inform. Transmission},
9(3):265--266, 1973.

\bibitem{LSTconst}
Nutan Limaye, Srikanth Srinivasan, and Sébastien Tavenas. Superpolynomial lower bounds
against low-depth algebraic circuits. \emph{Journal of the ACM}, 72(4):26:1--26:35, 2025.

\bibitem{MahajanVinay}
Meena Mahajan and V. Vinay. Determinant: combinatorics, algorithms, and complexity. \emph{Chicago
Journal of Theoretical Computer Science}, 1997(5), 1997.

\bibitem{MarkovShi}
Igor L. Markov and Yaoyun Shi. Simulating quantum computation by contracting tensor
networks. \emph{SIAM Journal on Computing}, 38(3):963--981, 2008.

\bibitem{MignonRessayre}
Thierry Mignon and Nicolas Ressayre. A quadratic bound for the determinant and permanent
problem. \emph{International Mathematics Research Notices}, 2004(79):4241--4253, 2004.

\bibitem{MulmuleySohoni}
Ketan D. Mulmuley and Milind Sohoni. Geometric complexity theory I: An approach to the P
vs. NP and related problems. \emph{SIAM Journal on Computing}, 31(2):496--526, 2001.

\bibitem{MulmuleySohoni2008}
Ketan D. Mulmuley and Milind Sohoni. Geometric complexity theory II: Towards explicit
obstructions for embeddings among class varieties. \emph{SIAM Journal on Computing}, 38(3):1175--1206, 2008.

\bibitem{Nisan}
Noam Nisan. Lower bounds for non-commutative computation. In \emph{Proceedings of the 23rd
Annual ACM Symposium on Theory of Computing}, pages 410--418, 1991.

\bibitem{NW97pd}
Noam Nisan and Avi Wigderson. Lower bounds on arithmetic circuits via partial derivatives.
\emph{Computational Complexity}, 6(3):217--234, 1997.

% The published bibliography skips numbers 57 and 58.
\setcounter{enumiv}{58}
\bibitem{CubOAI}
OpenAI. A cubic lower bound for border determinantal complexity of the permanent. OpenAI
Math Release preprint, September 24, 2026.
OAI:A-cubic-lower-bound-for-border-determinantal-complexity-of-the-permanent-September-24-2026.

\bibitem{DEP5OAI}
OpenAI. Homogeneous depth-five lower bounds for iterated matrix multiplication.
Preprint, September 25, 2026.
OAI:Homogeneous-depth-five-lower-bounds-for-iterated-matrix-multiplication-September-25-2026.

\bibitem{TenAdvances}
OpenAI. Ten Advances in Mathematics and Theoretical Computer Science. Preprint, updated
August 6, 2026.

\bibitem{ETHfalse}
OpenAI. The Exponential Time Hypothesis is False. OpenAI Math Release preprint, October
25, 2026. OAI:The-Exponential-Time-Hypothesis-is-False-October-25-2026.

\bibitem{PVCC}
Franco P. Preparata and Jean Vuillemin. The cube-connected cycles: a versatile network for
parallel computation. \emph{Commun. ACM}, 24(5):300--309, 1981.

\bibitem{Raghavan}
Prabhakar Raghavan. Probabilistic construction of deterministic algorithms: approximating
packing integer programs. \emph{J. Comput. System Sci.}, 37(2):130--143, 1988.

\bibitem{Raz06}
Ran Raz. Separation of multilinear circuit and formula size. \emph{Theory of Computing}, 2:121--135,
2006.

\bibitem{RazMultilinear}
Ran Raz. Multi-linear formulas for permanent and determinant are of super-polynomial size.
\emph{Journal of the ACM}, 56(2):8:1--8:17, 2009.

\bibitem{Raz13rank}
Ran Raz. Tensor-rank and lower bounds for arithmetic formulas. \emph{Journal of the ACM}, 60(6):40:1--40:15, 2013.

\bibitem{RSY08}
Ran Raz, Amir Shpilka, and Amir Yehudayoff. A lower bound for the size of syntactically
multilinear arithmetic circuits. \emph{SIAM Journal on Computing}, 38(4):1624--1647, 2008.

\bibitem{Sheshadri287}
Karthik Sheshadri. Extending OpenAI's cubic lower bound for the permanent to odd characteristic,
Pfaffian representations, and polynomial entries. ECCC Report TR26-287, 2026.

\bibitem{SheshadriBorder2026}
Karthik Sheshadri. A near-quadratic lower bound on the border determinantal complexity of
$\sum_i x_i^n$ via conormal specialization. arXiv:2606.13628v1, 2026.

\bibitem{ShpilkaYehudayoff}
Amir Shpilka and Amir Yehudayoff. Arithmetic circuits: A survey of recent results and open
questions. \emph{Foundations and Trends in Theoretical Computer Science}, 5(3--4):207--388, 2010.

\bibitem{StoehrVoloch}
Karl-Otto Stöhr and José Felipe Voloch. A formula for the Cartier operator on plane algebraic
curves. \emph{Journal für die reine und angewandte Mathematik}, 377:49--64, 1987.

\bibitem{StrassenBezout}
Volker Strassen. Die Berechnungskomplexität von elementarsymmetrischen Funktionen und von
Interpolationskoeffizienten. \emph{Numerische Mathematik}, 20:238--251, 1973.

\bibitem{Tavenas2013}
Sébastien Tavenas. Improved bounds for reduction to depth 4 and depth 3. \emph{Information and
Computation}, 240:2--11, 2015.

\bibitem{Uhlig}
Dietmar Uhlig. Networks computing Boolean functions for multiple input values. In M. S.
Paterson, editor, \emph{Boolean Function Complexity}, volume 169 of \emph{London Math. Soc. Lecture Note
Ser.}, pages 165--173. Cambridge University Press, 1992.

\bibitem{Valiant}
Leslie G. Valiant. Completeness classes in algebra. In \emph{Proceedings of the Eleventh Annual ACM
Symposium on Theory of Computing}, pages 249--261. ACM, 1979.

\bibitem{Val80}
Leslie G. Valiant. Negation can be exponentially powerful. \emph{Theoretical Computer Science},
12:303--314, 1980.

\bibitem{ValiantProjections}
Leslie G. Valiant. Reducibility by algebraic projections. \emph{L'Enseignement Mathématique}, 28:253--268, 1982.

\bibitem{VSBR}
Leslie G. Valiant, Sven Skyum, Stuart Berkowitz, and Charles Rackoff. Fast parallel computation
of polynomials using few processors. \emph{SIAM Journal on Computing}, 12(4):641--644, 1983.

\bibitem{AviICM}
Avi Wigderson. P, NP and mathematics---a computational complexity perspective. In \emph{Proceedings
of the ICM 2006 (Madrid)}, volume 1, pages 665--712. EMS, 2007.

\bibitem{Wbatch}
Richard Ryan Williams. Strong ETH Breaks With Merlin and Arthur: Short Non-Interactive
Proofs of Batch Evaluation. In \emph{Computational Complexity Conference}, volume 50 of \emph{LIPIcs},
pages 2:1--2:17, 2016.

\bibitem{Williams13}
Ryan Williams. Improving exhaustive search implies superpolynomial lower bounds. \emph{SIAM J.
Comput.}, 42:1218--1244, 2013.
\end{thebibliography}
